
\chapeter{Introducción}

La anemia falciforme (AF) es una hemoglobinopatía hereditaria con enorme impacto sanitario global. Estimaciones recientes del Global Burden of Disease indican que las personas que viven con AF aumentaron $\sim$41\% entre 2000 y 2021, de 5.46 millones a 7.74 millones; la mayor carga se concentra en África subsahariana, con alta mortalidad infantil en ausencia de diagnóstico y manejo oportunos. La OMS actualizó en 2025 sus lineamientos y fichas técnicas: subraya la subestimación histórica de muertes por AF y la necesidad crítica de expandir tamizaje y atención en entornos de bajos recursos.

Las vías diagnósticas estándar (electroforesis de hemoglobina, HPLC, pruebas confirmatorias específicas) son precisas, pero requieren infraestructura, reactivos y personal especializado; ello produce cuellos de botella logísticos y económicos que retrasan el diagnóstico temprano, sobre todo fuera de laboratorios de referencia. Esta brecha impulsa el desarrollo de diagnóstico en el punto de atención y estrategias portátiles de bajo costo capaces de tamizar con alta sensibilidad y razonable especificidad antes de confirmar en laboratorio.

En este contexto, la microscopía holográfica digital sin lente (\textit{lensless digital holographic microscopy}, DLHM) ofrece una alternativa hardware-ligera con gran ancho de banda espacio-frecuencia: el sensor registra hologramas en línea (patrones de interferencia entre la onda no desviada y la difractada por el objeto). El poder de DLHM reside en su simplicidad óptica, el gran campo visual y la posibilidad de construcción de bajo costo: existen prototipos open-source que integran fuente láser/LED, diafragma, carcasa 3D impresa, cámara tipo Raspberry y cómputo embarcado por menos de \$200, empleando algoritmos clásicos de Fresnel-Kirchhoff para reconstrucción digital cuando se requiere.

Tradicionalmente, el flujo holográfico reconstruye numéricamente la imagen (amplitud/fase) y luego aplica segmentación y clasificación. Sin embargo, la reconstrucción añade coste computacional y puede introducir artefactos (imagen gemela, aliasing, desenfoques). Una vía emergente es clasificar directamente los hologramas crudos, extrayendo de ellos descriptores discriminativos (textura, frecuencia, anisotropías) o bien usando transfer learning sobre el patrón de interferencia. Estudios recientes en holografía fuera de eje y en on-chip cytometry han demostrado que es factible alcanzar alta exactitud directamente sobre hologramas, reduciendo latencia e instrumentación computacional.

Desde el punto de vista físico, el holograma captura la huella interferométrica de la morfología eritrocitaria: un eritrocito sano (disco bicóncavo $\sim$8 $\mu$m) genera franjas aproximadamente concéntricas y de simetría casi circular; un drepanocito rígido/elongado (morfología de hoz) induce patrones más anisotrópicos con frecuencias espaciales dominantes más altas (franjas más juntas), dependientes además de la orientación. Estas diferencias se manifiestan en el dominio de Fourier (anillos/picos radiales y anisotropías angulares) y en texturas locales (microcontraste periódico), proporcionando señales robustas que no requieren la reconstrucción explícita para ser explotadas por algoritmos de clasificación. La literatura de DLHM detalla cómo la resolución lateral efectiva obedece a $\delta \approx \lambda z / L$ (longitud de onda $\lambda$, distancia objeto-sensor $z$, tamaño efectivo del holograma $L$), lo que enfatiza compromisos de muestreo y SNR; aun así, configuraciones compactas bien calibradas logran resoluciones micrométricas suficientes para tamizaje morfológico.

En inteligencia artificial para salud, tamaños muestrales modestos y heterogeneidad de adquisición exigen diseños estadísticos cuidadosos. La evidencia metodológica es clara: la validación cruzada simple puede sobreestimar el rendimiento con pocas muestras; en su lugar se recomienda validación cruzada anidada para desacoplar selección de hiperparámetros y estimación de desempeño, junto con bootstrap para intervalos de confianza. Estas prácticas reducen el sesgo optimista y mejoran la reproducibilidad, especialmente cuando se comparan modelos (SVM, bosques aleatorios, redes ligeras) y combinaciones de features (LBP/GLCM/FFT).

\textbf{Tesis del trabajo.} Este documento sostiene que es viable implementar un sistema de tamizaje portátil y de bajo costo para detección de anemia falciforme a partir de hologramas DLHM crudos de eritrocitos, sin reconstrucción explícita, usando un pipeline de preprocesamiento $\rightarrow$ extracción de rasgos físicos/estadísticos $\rightarrow$ clasificación interpretable entrenado y evaluado con protocolos estadísticos rigurosos. Las contribuciones son:

\begin{enumerate}
\item \textbf{Formulación físico-informacional de la tarea:} qué magnitudes del holograma (radiales/angulares en Fourier, texturas LBP/GLCM, medidas de anisotropía) correlacionan con morfología falciforme desde óptica de difracción.

\item \textbf{Ingeniería de características compacta y robusta al domain shift leve} (iluminación/contraste), incluyendo normalización y CLAHE.

\item \textbf{Comparativa de clasificadores clásicos y ligeros,} priorizando interpretabilidad y eficiencia computacional en hardware modestos.

\item \textbf{Evaluación con validez externa:} validación anidada, estimadores con IC95\% y análisis de errores por subgrupos (iluminación, artefactos).

\item \textbf{Puente hacia POCT:} discusión de integración con plataformas DLHM open-source y lineamientos de despliegue en campo.
\end{enumerate}

\textbf{Estructura del informe.} Tras esta introducción, se presenta el \textit{Planteamiento del problema} y su justificación clínica/tecnológica; un \textit{Marco teórico} que profundiza en holografía en línea, muestreo de Fresnel, ancho de banda espacio-frecuencia, y descriptores (LBP/GLCM/FFT) con formulación matemática; la \textit{Metodología} (datos, pipeline, validación, métricas e interpretabilidad); el \textit{Desarrollo} (implementación y decisiones de ingeniería), \textit{Resultados} (métricas con IC, ROC/PR, SHAP) y \textit{Conclusiones} (limitaciones, impacto y trabajo futuro). Las referencias se presentan en formato IEEE y privilegian revisiones y artículos recientes (2018-2025) en DLHM, IA biomédica y AF/POCT.